% =============================================================================
\section{Sensitivity to Hyperparameters}\label{app:sensitivity}
% =============================================================================

This appendix describes systematic ablation studies to assess ECL estimator sensitivity.

\subsection{Number of Monte Carlo samples $n$}

The convergence rate of $\widehat{\cI}(i;r)$ is $O(1/\sqrt{n})$ (\cref{thm:clt}). We recommend the following convergence check:
\begin{enumerate}[leftmargin=2em, label=(\roman*)]
\item Compute $\widehat{\ECL}_{0.9}$ at $n \in \{10, 25, 50, 100, 200, 500\}$.
\item Plot $\widehat{\ECL}_{0.9}$ vs.\ $n$ with bootstrap CIs.
\item Report the $n$ at which the CI width stabilizes below a target (e.g., $\pm$2~kb).
\end{enumerate}
Expected: ECL stabilizes rapidly for models with concentrated influence (small ECD) and more slowly for models with diffuse influence. The convergence assumption is falsified if $\Var(\widehat{\ECL}_{0.9} \mid n)$ does not decrease as $O(1/n)$ past $n \approx 100$, which would indicate either a non-CLT regime (e.g., heavy-tailed influence variables) or an estimator-specific bias not captured by the asymptotic theory.

\subsection{Block size $b$}

For block perturbation at sizes $b \in \{1, 8, 32, 128, 512\}$ bp:
\begin{enumerate}[leftmargin=2em, label=(\roman*)]
\item Compute ECL at each block size.
\item Report the block-size dependence curve.
\item Identify the coarsest block size that preserves ECL estimates within 10\% of single-nucleotide resolution.
\end{enumerate}
Expected: for models with smooth influence profiles, $b=32$ or $b=128$ suffices. For models with sharp motif-level features, $b=1$ may be necessary.

\subsection{Number of bootstrap replicates $B$}

Bootstrap CI coverage depends on $B$. We recommend $B \ge 1{,}000$ and verify coverage by checking that the CI width changes by less than 5\% when doubling $B$.

\subsection{Choice of embedding layer}

For multi-layer models, ECL can be computed at different layers $l$:
\begin{enumerate}[leftmargin=2em, label=(\roman*)]
\item Compute ECL at layers $l \in \{1, L/4, L/2, 3L/4, L\}$.
\item Expected: ECL grows with depth as deeper layers integrate more context.
\item Report the layer at which ECL saturates; this identifies the depth at which context integration is complete.
\end{enumerate}

\subsection{Choice of distance metric}

Compare ECL under: (i) squared Euclidean $\norm{\cdot}_2^2$, (ii) cosine distance, (iii) Mahalanobis distance with regularized covariance ($\lambda = 0.01$).

Expected: cosine distance yields shorter ECL than Euclidean (it ignores magnitude changes, which may extend further than directional changes). Mahalanobis distance should be intermediate, providing rotation-invariant measurement.

\subsection{Sensitivity summary}

\Cref{tab:hyperparameter_sensitivity} reports the results of this ablation study on a representative surrogate model (Enformer calibration, promoter locus). ECL$_{0.9}$ is stable across MC sample counts ($n \ge 25$), bootstrap replicates ($B \ge 500$), and embedding metric choice. Block size is the most influential parameter: coarsening from $b=1$ to $b=128$~bp reduces ECL$_{0.9}$ by $\sim$14\%, consistent with the smoothing effect discussed in \cref{alg:multiscale}.

\begin{table}[t]
\centering
\caption{Hyperparameter sensitivity: $\mathrm{ECL}_{0.9}$ and CI width under varying settings, on a representative locus drawn from the surrogate calibration of \cref{sec:experiments}. Estimates are stable across MC sample size, bootstrap replicates, and embedding metric, and degrade gracefully as block size grows.}
\label{tab:hyperparameter_sensitivity}
\begin{tabular}{llrr}
\toprule
\textbf{Hyperparameter} & \textbf{Value} & $\mathrm{ECL}_{0.9}$ (bp) & CI width (bp) \\
\midrule
MC samples (n) & 10 & 440 & 1 \\
MC samples (n) & 25 & 439 & 2 \\
MC samples (n) & 50 & 440 & 1 \\
MC samples (n) & 100 & 439 & 1 \\
MC samples (n) & 200 & 439 & 1 \\
\midrule
Block size (b) & 1 bp & 440 & n/a \\
Block size (b) & 8 bp & 404 & n/a \\
Block size (b) & 32 bp & 385 & n/a \\
Block size (b) & 128 bp & 378 & n/a \\
\midrule
Bootstrap (B) & 100 & 439 & 1 \\
Bootstrap (B) & 500 & 439 & 1 \\
Bootstrap (B) & 1000 & 439 & 1 \\
Bootstrap (B) & 5000 & 439 & 1 \\
\midrule
Metric & Squared Euclidean & 439 & n/a \\
Metric & Cosine & 440 & n/a \\
\bottomrule
\end{tabular}
\end{table}

